\documentclass[11pt,a4paper]{article}
\usepackage[utf8]{inputenc}
\usepackage[T1]{fontenc}
\usepackage[brazilian]{babel}
\usepackage[margin=2cm]{geometry}
\usepackage{lmodern}
\usepackage{amssymb}
\usepackage{xcolor}
\usepackage{array}
\usepackage{booktabs}
\usepackage{enumitem}
\usepackage{titlesec}
\usepackage{fancyhdr}
\usepackage{parskip}

\definecolor{brandgreen}{HTML}{2E5E3E}
\definecolor{lightbg}{HTML}{F5F2EC}

\titleformat{\section}{\large\bfseries\color{brandgreen}}{}{0pt}{}[{\color{brandgreen}\titlerule}]
\titleformat{\subsection}{\bfseries\color{brandgreen}}{}{0pt}{}

\pagestyle{fancy}
\fancyhf{}
\fancyhead[L]{\small Querência ERP --- Avaliação de usabilidade}
\fancyhead[R]{\small Fluxo 2 --- Associado}
\fancyfoot[C]{\small\thepage}
\renewcommand{\headrulewidth}{0.4pt}

\newcommand{\dizer}[1]{%
  \noindent\colorbox{lightbg}{%
    \begin{minipage}{0.95\linewidth}
    \vspace{1mm}
    \textbf{O que dizer ao participante:} ``#1''
    \vspace{1mm}
    \end{minipage}%
  }\vspace{3mm}
}
\newcommand{\criterio}[1]{\noindent\textit{Critério de sucesso objetivo:} #1}
\newcommand{\estado}[1]{\noindent\textit{Estado necessário (já preparado no kit):} #1\vspace{2mm}}

\begin{document}

\begin{center}
  {\LARGE\bfseries\color{brandgreen} Fluxo 2 --- Compra de ingresso}\\[1mm]
  {\large RF12, RF20, RF21 --- Papel: Associado (app mobile)}
\end{center}

\vspace{2mm}
\estado{conta de teste do kit, já logada pelo moderador; evento publicado com ingresso avulso
configurado.}

\begin{enumerate}[leftmargin=6mm]
  \item Abrir a aba \textbf{Eventos}, tocar no evento indicado no kit.
  \item Na tela do evento, tocar \textbf{Comprar ingressos}.
  \item Se aparecer o seletor ``Mesa / Ingresso avulso'', escolher \textbf{Ingresso avulso} (a
    menos que o moderador peça explicitamente a variante de reserva de mesa).
  \item Confirmar a compra do próprio ingresso (e, opcionalmente, adicionar 1 acompanhante).
  \item Abrir a aba \textbf{Ingressos}, localizar o ingresso recém-comprado --- confirmar que
    aparece um código (\texttt{\#001}, \texttt{\#002}\ldots) junto ao nome, e o QR ao tocar no
    card.
\end{enumerate}

\textbf{Verificação objetiva:}
\begin{small}
\begin{verbatim}
cd backend
node scripts/verificar-tarefas.js 2 <CPF-DA-CONTA-DE-TESTE> \
  "<trecho do nome do evento>"
\end{verbatim}
\end{small}
\criterio{ingresso emitido; check-in ainda não validado (normal nesse ponto).}

\vspace{2mm}
\noindent\textit{Nesta tarefa o QR do participante não é lido por ninguém --- a validação é tarefa
separada do participante da Diretoria, com ingressos à parte já pré-emitidos pelo moderador.}

\vfill
\noindent\footnotesize Se a tarefa travar por mais de \textasciitilde2 min, aplicar 1 intervenção
pontual e registrar; se mesmo assim não avançar, parar e marcar como \textit{não concluída} ---
protocolo completo no \texttt{roteiro-moderador.pdf}.

\end{document}
